\documentclass[border=10pt]{standalone}
\usepackage{tikz}
\usetikzlibrary{arrows.meta, shapes, calc}

\begin{document}
\begin{tikzpicture}[
    node distance=1.5cm,
    every node/.style={draw, thick, rounded corners, text width=4cm, align=center, inner sep=10pt, font=\sffamily},
    every path/.style={->, >=Stealth, blue, line width=1pt},
    base/.style={fill=white},
    theory/.style={base, fill=red!30},
    projection/.style={base, fill=green!50},
    algorithm/.style={base, fill=blue!30},
    result/.style={base, fill=orange!80},
    estimate/.style={base, fill=blue!50, text width=3cm},
    final/.style={base, fill=orange!50, text width=4cm}
]

% Top level
\node[theory] (cs) {Compressed Sensing Theory};

% Branches from top
\node[projection, below left=of cs] (pm) {Projection Matrix Theory};
\node[algorithm, below right=of cs] (ra) {Reconstruction Algorithms};

% Projection Matrix Theory children
\node[projection, below left=of pm] (opm) {Optimise Projection Matrix};
\node[projection, below right=of pm] (rnm) {Reduce Number of Measurements};

% Reconstruction Algorithms children
\node[algorithm, below left=of ra] (cr) {Convex Relaxation};
\node[algorithm, below=of cr] (s1) {Sparse Signal Estimate};

\node[algorithm, below right=of ra] (gp) {Greedy Pursuits};
\node[algorithm, below=of gp] (s2) {Sparse Signal Estimate};

\node[algorithm, below=of ra, right=of cr] (nc) {Non-Convex Minimisation Methods};
\node[algorithm, below=of nc] (s3) {Sparse Signal Estimate};

\node[algorithm, below=of ra, right=of gp] (ca) {Combinatorial Algorithms};
\node[algorithm, below=of ca] (s4) {Sparse Signal Estimate};

% Our Work
\node[result, below=of s1.south -| s2.south] (ow) {Our Work};

% Final result
\node[final, below=of ow] (f) {Improved Sparse Signal Estimate!};

% Connections
\draw (cs) -- (pm);
\draw (cs) -- (ra);

\draw (pm) -- (opm);
\draw (pm) -- (rnm);

\draw (ra) -- (cr);
\draw (ra) -- (gp);
\draw (ra) -- (nc);
\draw (ra) -- (ca);

\draw (cr) -- (s1);
\draw (gp) -- (s2);
\draw (nc) -- (s3);
\draw (ca) -- (s4);

\draw (s1) -- (ow);
\draw (s2) -- (ow);
\draw (s3) -- (ow);
\draw (s4) -- (ow);

\draw (ow) -- (f);

\end{tikzpicture}
\end{document}